% ===========================================================================
%  Raport — Faza 1: Cuantizare Statică FP16
%  Evaluare ViT-Tiny pretrained pe ImageNet-1k validation
%  Tudose Alexandru · Lucrare de Licență 2026
% ===========================================================================
\documentclass[12pt,a4paper]{article}

% --- Encoding & fonts (XeTeX / tectonic) ---
\usepackage{fontspec}
\defaultfontfeatures{Ligatures=TeX}

% --- Language ---
\usepackage{polyglossia}
\setmainlanguage{romanian}

% --- Page geometry ---
\usepackage[
  top=2.5cm, bottom=2.5cm,
  left=2.5cm, right=2.5cm
]{geometry}

% --- Math ---
\usepackage{amsmath, amssymb}

% --- Graphics ---
\usepackage{graphicx}
\usepackage{float}
\usepackage{caption}
\usepackage{subcaption}
\captionsetup{font=small, labelfont=bf}

% --- Tables ---
\usepackage{booktabs}
\usepackage{array}
\usepackage{multirow}

% --- Colors ---
\usepackage{xcolor}
\definecolor{fp32color}{HTML}{0072B2}
\definecolor{fp16color}{HTML}{009E73}
\definecolor{int8color}{HTML}{E69F00}
\definecolor{lightgray}{HTML}{F4F4F4}

% --- Hyperlinks ---
\usepackage[
  colorlinks=true,
  linkcolor=fp32color,
  citecolor=fp32color,
  urlcolor=fp32color
]{hyperref}

% --- Misc ---
\usepackage{microtype}
\usepackage{parskip}
\setlength{\parindent}{0pt}
\setlength{\parskip}{6pt}
\usepackage{enumitem}

% --- Header / Footer ---
\usepackage{fancyhdr}
\pagestyle{fancy}
\fancyhf{}
\lhead{\small Tudose Alexandru --- Faza 1: FP16 Static Quantization}
\rhead{\small 2026}
\cfoot{\thepage}
\renewcommand{\headrulewidth}{0.4pt}

% ===========================================================================
\begin{document}
% ===========================================================================

% --- Title page ---
\begin{titlepage}
  \centering
  {\large Universitatea Transilvania din Brașov\\}
  {\small Facultatea de Matematică și Informatică\\}
  {\small Specializarea: Informatică Aplicată\\[2cm]}

  \rule{\linewidth}{0.4pt}\\[0.5cm]
  {\LARGE\bfseries Faza 1: Cuantizare Statică FP16\\[0.3cm]}
  {\Large\bfseries Evaluare ViT-Tiny pretrained\\[0.2cm]}
  {\Large\bfseries pe ImageNet-1k validation\\[0.2cm]}
  {\large Post-Training Quantization --- FP32 vs.\ FP16}\\[0.3cm]
  \rule{\linewidth}{0.4pt}\\[2cm]

  \begin{minipage}{0.45\textwidth}
    \textbf{Student:}\\
    Tudose Alexandru\\
    Grupa 10LF333\\
    An III, Informatică Aplicată
  \end{minipage}
  \hfill
  \begin{minipage}{0.45\textwidth}
    \raggedleft
    \textbf{Coordonator:}\\
    Conf.\ dr.\ ing.\ Honorius Gâlmeanu\\
    Informatică Aplicată
  \end{minipage}\\[3cm]

\end{titlepage}

% --- TOC ---
\tableofcontents
\newpage

% ===========================================================================
\section{Introducere}
% ===========================================================================

Acest raport documentează \textbf{Faza~1} a lucrării de licență, în care
evaluăm impactul cuantizării statice FP16 asupra unui Vision Transformer
pretrained. Spre deosebire de studiul preliminar (Faza~0), care a antrenat
modelul pe CIFAR-10, această fază evaluează cuantizarea \emph{post-training}
--- fără nicio re-antrenare --- pe setul de validare ImageNet-1k.

\textbf{Obiective:}
\begin{enumerate}
  \item Stabilirea unui \textbf{baseline FP32} pe ImageNet-1k validation
    (50\,000 imagini, 1\,000 clase).
  \item Evaluarea conversiei \texttt{model.half()} (FP16 static) din punctul
    de vedere al acurateții, latenței și memoriei.
  \item Cuantificarea tradeoff-ului precizie--eficiență pentru FP16 pe un
    model pretrained real.
\end{enumerate}

\textbf{Context:} Cuantizarea FP16 este cea mai simplă formă de reducere a
preciziei numerice. Conversia \texttt{model.half()} transformă toate ponderile
din \texttt{float32} (4 bytes/param) în \texttt{float16} (2 bytes/param),
reducând teoretic memoria la \textbf{exact jumătate}. Pe hardware modern
(GPU-uri NVIDIA cu Tensor Cores, Apple Silicon cu suport nativ float16),
operațiile FP16 pot fi și mai rapide decât FP32.

% ===========================================================================
\section{Metodologie}
% ===========================================================================

\subsection{Model și date}

\begin{table}[H]
  \centering
  \caption{Configurația experimentală}
  \begin{tabular}{ll}
    \toprule
    \textbf{Parametru} & \textbf{Valoare} \\
    \midrule
    Model             & \texttt{vit\_tiny\_patch16\_224} (timm) \\
    Parametri         & 5.7 milioane \\
    Pretrained        & ImageNet-1k (\texttt{pretrained=True}) \\
    Dataset evaluare  & ImageNet-1k validation (50\,000 imagini, 1\,000 clase) \\
    Batch size        & 64 \\
    Device            & Apple M4, backend MPS \\
    PyTorch           & 2.9.1 \\
    \bottomrule
  \end{tabular}
  \label{tab:config_fp16}
\end{table}

\subsection{Preprocessing}

Transformarea de intrare este cea specificată de modelul timm:

\begin{itemize}
  \item Resize + center crop la $224 \times 224$ (crop ratio = 0.9)
  \item Interpolare bicubică
  \item Normalizare: $\mu = (0.5, 0.5, 0.5)$, $\sigma = (0.5, 0.5, 0.5)$
\end{itemize}

\textbf{Notă:} \texttt{vit\_tiny\_patch16\_224} din timm nu folosește
normalizarea standard ImageNet ($\mu = (0.485, 0.456, 0.406)$), ci
$(0.5, 0.5, 0.5)$ --- acesta este un aspect important, deoarece utilizarea
normalizării greșite ar degrada semnificativ acuratețea.

\subsection{Protocol de evaluare}

\begin{enumerate}
  \item \textbf{FP32 baseline:} Modelul pretrained este încărcat și evaluat
    direct în precizie \texttt{float32}.
  \item \textbf{FP16:} Un al doilea model pretrained identic este convertit
    cu \texttt{model.half()}, care transformă \emph{toate} ponderile și
    buffer-urile la \texttt{float16}.
  \item Ambele modele sunt evaluate pe aceleași 50\,000 de imagini,
    cu aceleași transformări de intrare.
  \item \textbf{Warmup:} Primele 3 batch-uri sunt excluse din măsurătorile
    de latență, pentru a elimina overhead-ul de inițializare (compilare
    kernele MPS, cache-uri).
  \item \textbf{Sincronizare:} După fiecare batch se apelează
    \texttt{torch.mps.synchronize()} pentru a asigura măsurători de latență
    corecte (operațiile MPS sunt asincrone).
\end{enumerate}

\subsection{Metrici}

\begin{itemize}
  \item \textbf{Top-1 Accuracy (\%):} Procentul de imagini clasificate corect.
  \item \textbf{Average Loss:} Cross-entropy loss mediu pe toate batch-urile.
  \item \textbf{Latency (ms/batch):} Timp mediu de inferență per batch de 64
    imagini, exclusiv warmup.
  \item \textbf{Model Memory (MB):} Memoria ocupată de parametrii modelului,
    calculată ca $\sum p.\text{numel()} \times p.\text{element\_size()}$.
\end{itemize}

% ===========================================================================
\section{Rezultate}
% ===========================================================================

\subsection{Comparație globală}

\begin{table}[H]
  \centering
  \caption{Rezultate FP32 vs.\ FP16 pe ImageNet-1k validation}
  \begin{tabular}{lrrrr}
    \toprule
    \textbf{Format} & \textbf{Accuracy (\%)} & \textbf{Loss} &
    \textbf{Latency (ms)} & \textbf{Memory (MB)} \\
    \midrule
    FP32 (baseline) & 75.456 & 0.9420 & 164.96 & 21.81 \\
    FP16 (\texttt{model.half()}) & 75.452 & 0.9420 & 135.14 & 10.91 \\
    \midrule
    \textbf{Diferență} & $\mathbf{-0.004}$~pp & $<0.001$ &
    $\mathbf{1.22\times}$ speedup & $\mathbf{2.00\times}$ reducere \\
    \bottomrule
  \end{tabular}
  \label{tab:fp16_results}
\end{table}

\begin{figure}[H]
  \centering
  \includegraphics[width=\textwidth]{../results/FP16ImageNet/plots/00_summary.png}
  \caption{Overview Faza~1: acuratețe, memorie și latență --- FP32 vs.\ FP16.
    Degradarea de acuratețe este invizibilă ($-0.004$~pp), reducerea de memorie
    este exactă ($2.00\times$), iar speedup-ul pe MPS este de $1.22\times$.}
  \label{fig:fp16_summary}
\end{figure}

\subsection{Acuratețe}

\begin{figure}[H]
  \centering
  \includegraphics[width=0.7\textwidth]{../results/FP16ImageNet/plots/01_accuracy_comparison.png}
  \caption{Comparație acuratețe Top-1: FP32 (75.456\%) vs.\ FP16 (75.452\%).
    Diferența de $-0.004$~pp este sub pragul de variabilitate numerică și
    confirmă că \texttt{model.half()} nu introduce degradare practică
    a acurateții pe ImageNet-1k.}
  \label{fig:fp16_accuracy}
\end{figure}

Degradarea de \textbf{$-0.004$~pp} este neglijabilă. Aceasta reprezintă
o diferență de doar \textbf{2 imagini} din 50\,000 clasificate diferit.
Loss-ul mediu este practic identic ($0.9420$ vs.\ $0.9420$), confirmând
că distribuția softmax nu este alterată semnificativ de reducerea preciziei.

\subsection{Memorie}

\begin{figure}[H]
  \centering
  \includegraphics[width=0.7\textwidth]{../results/FP16ImageNet/plots/02_memory_comparison.png}
  \caption{Footprint memorie parametri: FP32 (21.81~MB) vs.\ FP16 (10.91~MB).
    Reducerea este exact $2.00\times$, conform teoriei (4~bytes/param $\to$
    2~bytes/param).}
  \label{fig:fp16_memory}
\end{figure}

Reducerea de memorie de $\mathbf{2.00\times}$ este \emph{exact} conformă cu
teoria:
\begin{equation}
  \frac{21.81\text{ MB}}{10.91\text{ MB}} = 2.000\times
  \qquad \text{(4 bytes/param} \to \text{2 bytes/param)}
\end{equation}

Pentru un model mai mare (e.g., ViT-Base cu 86M parametri), acest factor
de $2\times$ ar traduce într-o economie de $\approx 330$~MB --- semnificativ
pentru deployment pe dispozitive edge.

\subsection{Latență}

\begin{figure}[H]
  \centering
  \includegraphics[width=0.7\textwidth]{../results/FP16ImageNet/plots/03_latency_comparison.png}
  \caption{Latență de inferență: FP32 (164.96~ms/batch) vs.\ FP16
    (135.14~ms/batch). Speedup de $1.22\times$ pe Apple M4 MPS.}
  \label{fig:fp16_latency}
\end{figure}

Speedup-ul de $\mathbf{1.22\times}$ este consistent cu arhitectura Apple M4:
\begin{itemize}
  \item Apple Silicon are suport nativ pentru operații \texttt{float16},
    cu unități ALU dedicate.
  \item Reducerea de bandwidth (jumătate din datele transferate între
    memorie și compute) contribuie direct la speedup.
  \item Speedup-ul sub $2\times$ este așteptat, deoarece nu toate operațiile
    beneficiază de FP16 (softmax, layer norm pot rămâne în FP32 intern).
\end{itemize}

\subsection{Tradeoff acuratețe--latență}

\begin{figure}[H]
  \centering
  \includegraphics[width=0.7\textwidth]{../results/FP16ImageNet/plots/04_accuracy_latency_tradeoff.png}
  \caption{Scatter plot acuratețe vs.\ latență. FP16 se află în colțul
    ideal (acuratețe menținută, latență redusă).}
  \label{fig:fp16_tradeoff}
\end{figure}

Graficul de tradeoff confirmă că FP16 domină strict pe axa eficienței:
acuratețea este practic identică, iar latența este redusă cu 30~ms/batch.
Nu există niciun compromis practic pentru a adopta FP16 pe acest model.

% ===========================================================================
\section{Discuție}
% ===========================================================================

\subsection{Validarea ipotezei}

Rezultatele confirmă ipoteza formulată în studiul preliminar (Faza~0):
\textbf{ViT-Tiny este robust la reducerea preciziei de la FP32 la FP16.}
Degradarea de $-0.004$~pp este cu un ordin de mărime mai mică decât
variabilitatea inter-run tipică ($\approx 0.1$~pp).

\subsection{Comparație cu studiul preliminar}

\begin{table}[H]
  \centering
  \caption{Comparație: Faza~0 (CIFAR-10) vs.\ Faza~1 (ImageNet-1k)}
  \begin{tabular}{lrr}
    \toprule
    \textbf{Metrică} & \textbf{Faza 0 (CIFAR-10)} & \textbf{Faza 1 (ImageNet-1k)} \\
    \midrule
    Accuracy FP32         & 82.88\% (fine-tuned)   & 75.456\% (pretrained)  \\
    Degradare FP16        & $<0.5$~pp              & $-0.004$~pp            \\
    Reducere memorie      & $2.00\times$           & $2.00\times$           \\
    Dataset               & 10\,000 val images     & 50\,000 val images     \\
    Clase                 & 10                     & 1\,000                 \\
    \bottomrule
  \end{tabular}
  \label{tab:fp16_vs_prelim}
\end{table}

Faza~1 confirmă pe un dataset mai complex (1\,000 clase, 50\,000 imagini)
ceea ce Faza~0 a sugerat pe CIFAR-10: \textbf{conversia FP16 este sigură
și fără compromisuri practice.}

\subsection{Implicații pentru fazele ulterioare}

Acuratețea FP32 de \textbf{75.456\%} devine baseline-ul de referință pentru
toate fazele ulterioare:
\begin{itemize}
  \item Faza~2 va evalua cuantizarea INT8 (weight-only, per-tensor) față de
    acest baseline.
  \item Faza~3 va extinde la INT8 per-channel și va realiza analiza de
    sensibilitate per strat.
  \item Orice degradare $<1$~pp față de 75.456\% este considerată acceptabilă
    pentru o reducere de memorie semnificativă.
\end{itemize}

\subsection{Limitări}

\begin{enumerate}
  \item \textbf{Hardware-specific:} Speedup-ul de $1.22\times$ este specific
    Apple M4 MPS. Pe GPU-uri NVIDIA cu Tensor Cores, speedup-ul poate fi
    semnificativ mai mare (până la $2$--$3\times$).
  \item \textbf{Un singur model:} Experimentul evaluează doar ViT-Tiny.
    Modele mai mari (ViT-Base, ViT-Large) ar putea beneficia mai mult
    de reducerea de bandwidth.
  \item \textbf{Batch size fix:} Batch-ul de 64 este limitat de memoria
    disponibilă. Cu FP16, ar fi posibil un batch size mai mare, ceea ce
    ar putea îmbunătăți throughput-ul total.
  \item \textbf{Procese de fundal:} Măsurătorile de latență au fost
    efectuate cu alte procese active (browser, Discord), ceea ce poate
    introduce variabilitate. Totuși, raportul de speedup rămâne valid
    deoarece ambele configurații au fost evaluate în aceleași condiții.
\end{enumerate}

% ===========================================================================
\section{Concluzii}
% ===========================================================================

\begin{enumerate}
  \item \textbf{Degradare neglijabilă:} FP16 introduce o pierdere de
    acuratețe de doar $\mathbf{-0.004}$~pp (2 imagini din 50\,000),
    practic invizibilă.

  \item \textbf{Reducere de memorie exactă:} $\mathbf{2.00\times}$,
    conform cu teoria ($\texttt{float32} \to \texttt{float16}$:
    4~bytes $\to$ 2~bytes per parametru).

  \item \textbf{Speedup pe Apple M4:} $\mathbf{1.22\times}$ reducere
    a latenței de inferență, datorită suportului nativ float16 și
    reducerii de bandwidth.

  \item \textbf{Recomandare:} \texttt{model.half()} este o optimizare
    \emph{fără cost} pentru ViT-Tiny pe ImageNet-1k --- poate fi aplicată
    universal ca prim pas de optimizare, înaintea tehnicilor mai agresive
    (INT8, pruning, distillation).

  \item \textbf{Baseline stabilit:} Acuratețea FP32 de 75.456\% pe
    ImageNet-1k validation devine referința pentru evaluarea cuantizării
    INT8 în Faza~2.
\end{enumerate}

% ===========================================================================
\begin{thebibliography}{9}
% ===========================================================================

\bibitem{dosovitskiy2021vit}
A. Dosovitskiy, L. Beyer, A. Kolesnikov et al.,
``An Image is Worth 16x16 Words: Transformers for Image Recognition at Scale,''
\emph{ICLR}, 2021. \href{https://arxiv.org/abs/2010.11929}{arXiv:2010.11929}.

\bibitem{touvron2021deit}
H. Touvron, M. Cord, M. Douze et al.,
``Training data-efficient image transformers \& distillation through attention,''
\emph{ICML}, 2021. \href{https://arxiv.org/abs/2012.12877}{arXiv:2012.12877}.

\bibitem{jacob2018quantization}
B. Jacob, S. Kligys, B. Chen et al.,
``Quantization and Training of Neural Networks for Efficient Integer-Arithmetic-Only Inference,''
\emph{CVPR}, 2018. \href{https://arxiv.org/abs/1712.05877}{arXiv:1712.05877}.

\bibitem{timm}
R. Wightman, ``PyTorch Image Models (timm),'' 2019.
\url{https://github.com/rwightman/pytorch-image-models}.

\bibitem{imagenet}
O. Russakovsky, J. Deng, H. Su et al.,
``ImageNet Large Scale Visual Recognition Challenge,''
\emph{IJCV}, vol.~115, no.~3, pp.~211--252, 2015.
\href{https://arxiv.org/abs/1409.0575}{arXiv:1409.0575}.

\end{thebibliography}

\end{document}
